\documentclass[12pt]{article}
\usepackage[utf8]{inputenc}
\usepackage[spanish]{babel}
\usepackage[a4paper, margin=2cm]{geometry}
\usepackage{tikz}
\usetikzlibrary{babel}
\usepackage[most]{tcolorbox}
\usepackage{amsmath}

% Usar fuente sin serifa para un aspecto más moderno (tipo infografía)
\renewcommand{\familydefault}{\sfdefault}

% Definición de colores con temática médica/científica
\definecolor{medblue}{RGB}{43, 108, 176}
\definecolor{medteal}{RGB}{49, 151, 149}
\definecolor{medlight}{RGB}{235, 248, 255}
\definecolor{meddark}{RGB}{26, 32, 44}
\definecolor{mercurygray}{RGB}{156, 163, 175}  % Gris para el mercurio
\definecolor{monitorbg}{RGB}{17, 24, 39}       % Fondo oscuro de pantalla digital
\definecolor{monitorgreen}{RGB}{16, 185, 129}  % Verde de lectura digital

\begin{document}
\pagestyle{empty}

\begin{tcolorbox}[
    enhanced, 
    colback=medlight!30, 
    colframe=medblue, 
    title={\Large \textbf{Sistemas de Unidades: Presión Arterial}}, 
    halign title=center, 
    fonttitle=\bfseries, 
    arc=4mm, 
    boxrule=1.5pt,
    shadow={2mm}{-2mm}{0mm}{black!30!white}
]

    \vspace{0.2cm}
    \begin{tcolorbox}[colback=white, colframe=medteal, arc=2mm, boxrule=1.2pt, title=\textbf{Caso Clínico / Problema}]
        La presión sistólica normal es de unos \textbf{120 mmHg} (milímetros de mercurio). Conociendo que \textbf{$1 \text{ mmHg} \approx 133.322 \text{ Pa}$} (Pascales, medida del Sistema Internacional), convierta la presión sistólica a Pascales.
    \end{tcolorbox}

    \vspace{0.4cm}
    
    \begin{center}
    \begin{tikzpicture}[scale=0.95]
        
        % --- Columna de Mercurio (Izquierda) ---
        \begin{scope}[shift={(-3, 0)}]
            % Base del esfigmomanómetro
            \filldraw[meddark!10, draw=meddark, thick, rounded corners=2pt] (-1, -1.2) rectangle (0.8, 6);
            
            % Tubo de vidrio
            \filldraw[white, draw=meddark!50, thick] (-0.4, -0.5) rectangle (0.1, 5.5);
            
            % Bulbo y Mercurio
            \filldraw[mercurygray, draw=meddark, thick] (-0.4, -0.5) rectangle (0.1, 3);
            \filldraw[mercurygray, draw=meddark, thick] (-0.15, -0.6) circle (0.45cm);
            
            % Escala
            \foreach \y/\val in {0/0, 1/40, 2/80, 3/120, 4/160, 5/200} {
                \draw[thick, meddark] (0.1, \y) -- (0.3, \y) node[right, font=\scriptsize\bfseries] {\val};
            }
            
            % Indicador de 120
            \draw[<-, thick, red!80!black] (0.9, 3) -- (1.5, 3) node[right, font=\bfseries] {120 mmHg};
            
            \node[font=\bfseries, text=meddark, align=center] at (-0.1, -2.1) {Esfigmomanómetro\\(Columna de Hg)};
        \end{scope}

        % --- Flecha de conversión ---
        \draw[->, ultra thick, medteal] (0.3, 2.5) -- (2.6, 2.5) node[midway, above, font=\small\bfseries] {$\times 133.322$};
        \draw[->, ultra thick, medteal] (0.3, 2.5) -- (2.6, 2.5) node[midway, below, font=\small\bfseries] {Conversión al S.I.};

        % --- Monitor Digital (Derecha) ---
        \begin{scope}[shift={(5.5, 1.5)}]
            % Carcasa del monitor
            \filldraw[meddark, rounded corners=5pt] (-2.5, -1.5) rectangle (2.5, 3.5);
            % Pantalla
            \filldraw[monitorbg, draw=meddark!80, thick, rounded corners=3pt] (-2, -0.5) rectangle (2, 2.5);
            
            % Lectura digital
            \node[font=\huge\bfseries, text=monitorgreen] at (0, 1.4) {15 998};
            \node[font=\large\bfseries, text=monitorgreen!70] at (0, 0.3) {Pa};
            
            % Botones decorativos
            \filldraw[medteal] (-1.5, -1) circle (0.2cm);
            \filldraw[medteal] (-0.8, -1) circle (0.2cm);
            \filldraw[red!80] (1.5, -1) circle (0.2cm);
            
            \node[font=\bfseries, text=meddark, align=center] at (0, -2.6) {Monitor Digital\\(Pascales)};
        \end{scope}
        
    \end{tikzpicture}
    \end{center}

    \vspace{0.4cm}
    \begin{tcolorbox}[colback=medteal!10, colframe=medteal, arc=2mm, boxrule=1.2pt, title=\textbf{Desarrollo Matemático y Solución}]
        Para convertir la presión sistólica de milímetros de mercurio (mmHg) a Pascales (Pa), utilizamos el factor de conversión proporcionado:
        \begin{align*}
            P &= 120 \text{ mmHg} \times \left( \frac{133.322 \text{ Pa}}{1 \text{ mmHg}} \right) \\[1ex]
            P &= 120 \times 133.322 \text{ Pa} = 15\,998.64 \text{ Pa}
        \end{align*}
        Podemos redondear este valor y expresarlo en notación científica para mayor comodidad:
        \[ P \approx \mathbf{1.6 \times 10^4 \text{ Pa}} \]
        \textbf{Resultado:} La presión sistólica normal de $120$ mmHg equivale a \textbf{15,998.64 Pa} (aprox. $1.6 \times 10^4$ Pa).
    \end{tcolorbox}
    
\end{tcolorbox}

\end{document}